\heading{热学-模拟题1-期中}

\noteinfo{题目和答案由AI生成。}

\textbf{一、简答题}
\begin{question}
写出重力场中处于热平衡的稀薄理想气体数密度分布 $n(z)$，并求单个分子的平均势能。
\begin{solution}

Boltzmann 分布给出
\[
 n(z)=n_0 e^{-mgz/(k_BT)}.
\]
归一化后的概率密度可写为
\[
 f(z)=\frac{mg}{k_BT}e^{-mgz/(k_BT)},\qquad z>0.
\]
故平均势能
\[
\langle U\rangle=\int_0^{\infty}mgz\,f(z)\,dz=k_BT.
\]
\ansmath{\langle U\rangle=\int_0^{\infty}mgz\,f(z)\,dz=k_BT}

\end{solution}
\end{question}

\begin{question}
说明为什么等温大气中上升的理想气球所受浮力与高度无关。
\begin{solution}

等温大气中
\[
\rho(z)=\frac{p(z)\mu}{RT}.
\]
若气球中气体物质的量不变，则
\[
V(z)=\frac{p_0V_0}{p(z)}.
\]
故浮力
\[
F_b=\rho(z)gV(z)=\frac{p_0V_0\mu g}{RT},
\]
与 $z$ 无关。
\ansmath{F_b=\rho(z)gV(z)=\frac{p_0V_0\mu g}{RT},}

\end{solution}
\end{question}

\begin{question}
写出泻流数（单位时间单位面积穿过小孔的粒子数）和碰撞频率的公式。
\begin{solution}

\[
\Gamma=\frac14 n\langle v\rangle,
\qquad
\langle Z\rangle=\sqrt2\,n\sigma\langle v\rangle,
\]
其中 $\sigma=\pi d^2$ 为有效碰撞截面。
\ansmath{\Gamma=\frac14 n\langle v\rangle,
\qquad
\langle Z\rangle=\sqrt2\,n\sigma\langle v\rangle,}

\end{solution}
\end{question}

\begin{question}
利用能均分定理写出单原子、刚性双原子、非刚性双原子的平均内能。
\begin{solution}

\[
\langle\varepsilon\rangle=\frac32k_BT,\qquad \frac52k_BT,\qquad \frac72k_BT.
\]
\ansmath{\langle\varepsilon\rangle=\frac32k_BT,\qquad \frac52k_BT,\qquad \frac72k_BT}

\end{solution}
\end{question}

\begin{question}
为什么 Feynman 棘轮在与单一热源接触时长期平均不能输出净功？
\begin{solution}

因为棘轮与卡扣都受同一热源的热涨落控制，顺转与逆转都以相同的玻尔兹曼因子被激发，长期平均无净转动；否则将违背第二定律。
\anstext{见上。}

\end{solution}
\end{question}

\textbf{二、计算与证明题}
\begin{question}
从 Maxwell 速率分布出发，求速率大于 $v_0$ 的分子在单位时间单位面积上撞击容器壁的次数。
\begin{solution}

对指向容器壁的半空间积分，
\[
\Gamma(v_0)=\int_{v\ge v_0,\ v_x>0}v_x f(\mathbf v)\,d^3v.
\]
将角度积分做完，得到
\[
\Gamma(v_0)=\frac{n}{4}\sqrt{\frac{8k_BT}{\pi m}}
\left(1+\frac{mv_0^2}{2k_BT}\right)e^{-mv_0^2/(2k_BT)}.
\]
这是题中常用结论。
\ansmath{\Gamma(v_0)=\frac{n}{4}\sqrt{\frac{8k_BT}{\pi m}}
\left(1+\frac{mv_0^2}{2k_BT}\right)e^{-mv_0^2/(2k_BT)}}

\end{solution}
\end{question}

\begin{question}
设一竖直圆筒内气体满足 van der Waals 近似
\[
\left(p+\frac{a}{V_m^2}\right)(V_m-b)=RT.
\]
若忽略分子间引力项 $a$，证明其状态方程可写成
\[
\frac{p(V-b)}{T}=C.
\]

\begin{solution}

忽略 $a$ 后有
\[
p(V-b)=RT.
\]
若以固定物质的量写成总体积形式，右边只差常数因子，故直接整理得
\ansmath{\frac{p(V-b)}{T}=C.}
若从全微分形式出发，也可写为
\[
\frac{dp}{p}=\frac{dT}{T}-\frac{dV}{V-b},
\]
积分后即得同样结果。
\end{solution}
\end{question}

\begin{question}
设两个热容分别为 $C_1,C_2$ 的物体，初温分别为 $T_{10}>T_{20}$，通过一块热导率常数为 $\kappa$ 的薄板交换热量，满足
\[
\dot Q=\kappa(T_1-T_2).
\]
求温差减半所需时间。
\begin{solution}

能量守恒给出
\[
-C_1\dot T_1=C_2\dot T_2=\kappa(T_1-T_2).
\]
令 $\Delta=T_1-T_2$，则
\[
\dot\Delta=\dot T_1-\dot T_2=-\kappa\left(\frac1{C_1}+\frac1{C_2}\right)\Delta.
\]
故
\[
\Delta(t)=\Delta_0\exp\!\left[-\kappa\frac{C_1+C_2}{C_1C_2}t\right].
\]
令 $\Delta(t_{1/2})=\Delta_0/2$，得
\ansmath{t_{1/2}=\frac{C_1C_2}{\kappa(C_1+C_2)}\ln2.}
\end{solution}
\end{question}

\begin{question}
推导圆管层流的 Poiseuille 定律。已知液体粘滞系数为 $\eta$，管长 $L$、半径 $R$，压强差为 $\Delta p$。
\begin{solution}

取半径为 $r$ 的圆柱壳，受压强驱动力与粘滞阻力平衡：
\[
\Delta p\,\pi r^2=2\pi rL\,\eta\left(-\frac{du}{dr}\right).
\]
故
\[
\frac{du}{dr}=-\frac{\Delta p}{2\eta L}r.
\]
结合无滑移边界条件 $u(R)=0$，得速度分布
\[
 u(r)=\frac{\Delta p}{4\eta L}(R^2-r^2).
\]
流量
\[
Q=\int_0^R2\pi r u(r)\,dr=\frac{\pi R^4}{8\eta L}\Delta p.
\]
即
\ansmath{Q=\frac{\pi R^4}{8\eta L}\Delta p.}
\end{solution}
\end{question}

\begin{question}
半径为 $r$ 的小球悬浮在液体中，液体和小球密度分别为 $\rho_0,\rho$，在重力场中达到热平衡。测得高度差 $\Delta h$ 处的平衡数密度分别为 $n_1,n_2$。推导 Avogadro 常数的计算公式。
\begin{solution}

小球在高度 $z$ 处的有效重力势能为
\[
 U(z)=\frac43\pi r^3(\rho-\rho_0)gz.
\]
由 Boltzmann 分布
\[
 n(z)\propto e^{-N_AU(z)/(RT)}.
\]
因此
\[
\frac{n_1}{n_2}=\exp\left(\frac{4\pi r^3(\rho-\rho_0)g\Delta h}{3RT}N_A\right).
\]
解得
\ansmath{N_A=\frac{3RT}{4\pi r^3(\rho-\rho_0)g\Delta h}\ln\frac{n_1}{n_2}.}
\end{solution}
\end{question}

\begin{question}
设大气温度线性递减：
\[
T(z)=T_0-\alpha z,
\qquad 0\le z<T_0/\alpha.
\]
求压强分布 $p(z)$。
\begin{solution}

静力平衡
\[
\frac{dp}{dz}=-\rho g.
\]
理想气体状态方程给出
\[
\rho=\frac{p\mu}{RT(z)}=\frac{p\mu}{R(T_0-\alpha z)}.
\]
故
\[
\frac{dp}{p}=-\frac{\mu g}{R(T_0-\alpha z)}dz.
\]
积分得
\[
\ln\frac{p(z)}{p_0}=\frac{\mu g}{R\alpha}\ln\left(1-\frac{\alpha z}{T_0}\right).
\]
因此
\ansmath{p(z)=p_0\left(1-\frac{\alpha z}{T_0}\right)^{\mu g/(R\alpha)}.}
当 $\alpha\to0$ 时回到等温大气公式 $p(z)=p_0e^{-\mu gz/(RT_0)}$。
\end{solution}
\end{question}


